\documentclass{article}
\usepackage{mate}

\title{Serii}
\author{}
\date{}

\begin{document}

\maketitle

\section{Limite de siruri complexe}
  \begin{enumerate}

    \definitie Fie $(z_n)_{n \geq q}$ un sir de numere complexe. Sirul este \term{convergent} daca $\exists L \in \mathbb{C}$ astfel incat $\forall \varepsilon > 0, \exists n_{\varepsilon} \in \mathbb{N}$ astfel incat $\forall n \geq n_{\varepsilon}$ avem $\abs{z_n - L} < \varepsilon$. (Orice vecinatate a numarului complex $L$ contine toti termenii sirului, exceptand cel mult un numar finit)
    \definitie Un sir este \term{divergent} daca nu este convergent.

\notatie Daca exista limita, atunci aceasta este unica si se noteaza $L = \lim_{n \to \infty} z_n$ sau $z_n \to L$

\observatie $z_n \to L \Leftrightarrow \abs{z_n - L} \to 0$ .
\observatie $z_n \to 0 \Leftrightarrow \abs{z_n} \to 0$.
\observatie $z_n \to L \Rightarrow \abs{z_n} \to \abs{L}$.
\proprietate Daca $z_n \to L_1$ si $v_n \to L_2$, atunci $z_n + v_n \to L_1 + L_2$.
\proprietate Daca $z_n \to L_1$ si $v_n \to L_2$, atunci $z_n \cdot v_n \to L_1 \cdot L_2$.
\proprietate Daca $z_n \to L$ si $\alpha \in \mathbb{R}$, atunci $\alpha z_n \to \alpha L$
\proprietate Daca $z_n \to L_1$ si $v_n \to L_2 \neq 0$, atunci $\exists N \in \mathbb{N}$ astfel incat $ v_n \neq 0, \forall n \geq N$ si $\frac{z_n}{v_n} \to \frac{L_1}{L_2}$
\proprietate Daca $z_n = x_n + iy_n$, unde $x_n, y_n \in \mathbb{R}$, iar $L = a + ib$, unde $a, b \in \mathbb{R}$, avem $z_n \to L \Leftrightarrow x_n \to a \text{ si } y_n \to b$

\proprietate Daca un sir de numere reale este monoton si marginit, atunci este convergent.
\proprietate Daca un sir de numere reale este monoton si nemarginit, atunci este divergent.
\proprietate Orice sir de numere reale monoton are limita, finita sau infinita.

\propozitie (Criteriul majorarii) Fie $(z_n)_{n \geq q}$ un sir de numere complexe si $(a_n)_{n \geq N}, N \geq q$ un sir de numere reale cu $a_n \to 0$, astfel incat $\abs{z_n - L} \leq a_n, \forall n \geq N$, atunci $z_n \to L$. 
\observatie Daca un sir de numere complexe este convergent, atunci este si marginit ($\exists M \geq 0$ astfel incat $\abs{z_n} \leq M, \forall n \geq q)$.

\definitie Fie $(z_n)_{n \geq q}$ un sir de numere complexe si $(k_n)_{n \geq p}$ un sir strict crescator de numere naturale, unde $k_n \geq q, \forall n \geq p$. Sirul $(z_{k_n})_{n \geq p}$ se numeste \term{subsir} al sirului $(z_n)_{n \geq q}$.

\propozitie Fie $(z_n)_{n \geq q}$ un sir de numere complexe si $(z_{k_n})_{n \geq p}$ un subsir al sau. Daca $z_n \to L$, atunci $z_{k_n} \to L$. (Daca un sir este convergent, orice subsir al sau este convergent la aceeasi limita)
\proprietate Daca $z_n \to L$, atunci $z_{n + k} \to L, \forall k \geq 0$.

\exemplu $z_n = \frac{n}{2n+1} + i\left(1+\frac{1}{n}\right)^n \to \frac{1}{2} + ie$
\exemplu $z_n = (-1)^n + i\left(1 + \frac{1}{n}\right)^n$ este divergent.



\definitie Sirul $(z_n)_{n \geq q}$ se numeste sir Cauchy daca $\forall \varepsilon > 0, \exists n_{\varepsilon} \geq q$ astfel incat $\forall n \geq n_{\varepsilon}, p \in \mathbb{N}^*$ avem $\abs{z_{n+p} - z_n} < \varepsilon$. (Distanta dintre oricare doi termeni devine arbitrar de mica)

\propozitie Fie $(z_n)_{n \geq q}$ un sir de numere complexe. Atunci $(z_n)_{n \geq q}$ este convergent daca si numai daca este sir Cauchy.

\observatie Aceasta relatie nu are loc in orice spatiu metric. De exemplu, in $\mathbb{Q}$ exista siruri Cauchy care nu au limita rationala.

\section{Convergenta unei serii}

\notatie Seriile se noteaza cu forma $\sum_{n \geq q} z_n$.
\definitie Sirul $(S_n)_{n \geq q}$, $S_n = \sum_{k = q}^{n}z_{k}$ se numeste \term{sirul de sume partiale} asociat seriei $\sum_{n \geq q} z_n$.
\definitie Seria $\sum_{n \geq q} z_n$ este \term{convergenta} daca sirul sau de sume partiale este convergent.
\definitie Daca o serie $\sum_{n \geq q} z_n$ este convergenta, atunci \term{suma seriei} este \\ $\lim_{n \to \infty} S_n \xlongequal{\text{not}} \sum_{n = q}^{\infty} z_n \in \mathbb{C}$


\proprietate Fie $p, q \in \mathbb{N}, p > q$. Atunci seria $\sum_{n \geq q} z_n$ convergenta $\Leftrightarrow \sum_{n \geq p} z_n$ convergenta.

\proprietate Daca $\sum_{n \geq q} z_n$ si $\sum_{n \geq q} v_n$ sunt serii convergente, atunci seria $\sum_{n \geq q} (\alpha z_n + \beta v_n)$ este convergenta, $\forall \alpha, \beta \in \mathbb{C}$ iar suma $\sum_{n = q}^{\infty} (\alpha z_n + \beta v_n) = \alpha \sum_{n = q}^{\infty}z_n + \beta \sum_{n = q}^{\infty}v_n$. (Combinatia liniara de serii convergente este egala cu seria combinatiei liniare)

\proprietate Daca $\sum_{n \geq q} z_n$ este convergenta, iar $\sum_{n \geq q} v_n$ este divergenta, atunci $\sum_{n \geq q} (z_n + v_n)$ este divergenta.

\proprietate Daca $\sum_{n \geq q} z_n$ este o serie divergenta si $\alpha \in \mathbb{C}^*$, atunci seria $\sum_{n \geq q} \alpha z_n$ este divergenta.

\begin{tcolorbox}
\propozitie (Criteriul necesar) Daca $\sum_{n \geq q} z_n$ este convergenta, atunci $z_n \to 0$.
\end{tcolorbox}

\propozitie (Seria geometrica) Fie $z$ un numar complex. Daca $\abs{z} < 1$, atunci $\sum_{n \geq 0} z^n$ este convergenta si \\ $\sum_{n = 0}^{\infty} z^n = \frac{1}{1-z}$. Daca $\abs{z} \geq 1$, atunci $\sum_{n \geq 0} z^n$ este divergenta.

\observatie Daca $p_1, p_2, q \in \mathbb{N}, p_1 \neq p_2$ si $p_1, p_2 > q$, avem $\sum_{n \geq p_1} z_n$ convergenta $\Leftrightarrow \sum_{n \geq p_2} z_n$ convergenta.


\corolar Daca $\sum_{n \geq q} z_n$ este o serie convergenta si $\alpha \in \mathbb{C}$, atunci seria $\sum_{n \geq q} \alpha z_n$ este convergenta, iar suma ei $\sum_{n = q}^{\infty} \alpha z_n = \alpha \sum_{n = q}^{\infty} z_n $. (Scalarul iese in fata la serii convergente)




\propozitie (Criteriul Cauchy) Fie $(z_n)_{n \geq q}$ sir de nr complexe. Seria $\sum_{n \geq q} z_n$ convergenta daca si numai daca sirul sumelor partiale este sir Cauchy, adica $\forall \varepsilon > 0, \exists n_{\varepsilon} \geq q$ astfel incat $\forall n \geq n_{\varepsilon}, p \in \mathbb{N}^*$ avem $\abs{\sum_{k= n}^{n+p} z_k} < \varepsilon$.



\end{enumerate}
\section{Serii cu termeni pozitivi}
\begin{enumerate}

\observatie Daca o serie are doar termeni reali pozitivi, atunci sirul sumelor partiale este crescator, deoarece $S_{n+1} - S_n = x_{n+1} \geq 0, \forall n$.
\observatie O astfel de serie este convergenta daca si numai daca sirul sumelor partiale este marginit superior.

\begin{tcolorbox}
\definitie O serie de forma $\sum_{n \geq 1} \frac{1}{n^\alpha}$, $\alpha \in \mathbb{R}$ se numeste \term{serie Riemann}.
\observatie O astfel de serie este convergenta $\Leftrightarrow \alpha > 1$.
\end{tcolorbox}


\begin{tcolorbox}
\propozitie (Criteriul de comparatie cu inegalitati) Fie doua serii cu termeni pozitivi $\sum_{n \geq q} x_n$ si $\sum_{n \geq q} y_n$.
  \begin{enumerate}
    \item Daca $\exists N \geq q$ astfel incat $\forall n \geq N$ avem $0 \leq x_n \leq y_n$, iar $\sum_{n} y_n$ este convergenta, atunci $\sum_{n} x_n$ este convergenta.
    \item Daca $\exists N \geq q$ astfel incat $\forall n \geq N$ avem $x_n \geq y_n \geq 0$, iar $\sum_{n} y_n$ este divergenta, atunci $\sum_{n} x_n$ este divergenta.
  \end{enumerate}
\end{tcolorbox}

\begin{tcolorbox}
  \propozitie (Criteriul de comparatie la limita) Fie doua serii $\sum_{n \geq q}x_n$ si $\sum_{n \geq q}y_n$, $N \geq q$ astfel incat $\forall n \geq N, x_n \geq 0$ si $y_n > 0$. Presupunem ca $\exists \lim_{n \to \infty} \frac{x_n}{y_n} \xlongequal{\text{not}} L \in \overline{\mathbb{R}}$.

  \begin{enumerate}
    \item Daca $L$ este finit si $\sum_{n} y_n$ este convergenta, atunci $\sum_{n} x_n$ este convergenta.
    \item Daca $L \neq 0$ si $\sum_{n} y_n$ este divergenta, atunci $\sum_{n} x_n$ este divergenta.
    \item Daca $L$ este finit si nenul, atunci $\sum_{n} x_n$ si $\sum_{n} y_n$ au aceeasi natura.

  \end{enumerate}
\end{tcolorbox}

\begin{tcolorbox}
\propozitie (Criteriul radical) Fie $N > 0$ si $x_n \geq 0, \forall n \geq N$, astfel incat $\sqrt[n]{x_n} \to L \in \mathbb{R}$.

\begin{enumerate}
  \item Daca $L < 1$, atunci $\sum_{n} x_n$ este convergenta.
  \item Daca $L > 1$, atunci $\sum_{n} x_n$ este divergenta.
\end{enumerate}
\end{tcolorbox}

\begin{tcolorbox}
\propozitie (Criteriul raportului) Fie $N > 0$ si $x_n > 0, \forall n \geq N$, astfel incat \\ $\frac{x_{n+1}}{x_n} \to L \in \mathbb{R}$.

\begin{enumerate}
  \item Daca $L < 1$, atunci $\sum_{n} x_n$ este convergenta.
  \item Daca $L > 1$, atunci $\sum_{n} x_n$ este divergenta.
\end{enumerate}
\end{tcolorbox}

\begin{tcolorbox}
\propozitie (Criteriul Raabe-Duhamel) Fie $N > 0$ si $x_n > 0, \forall n \geq N$, astfel incat \\ $n \left(\frac{x_n}{x_{n+1}} - 1 \right) \to L \in \mathbb{R}$.

  \begin{enumerate}
    \item Daca $L < 1$, atunci $\sum_{n} x_n$ este divergenta (!).
    \item Daca $L > 1$, atunci $\sum_{n} x_n$ este convergenta (!).
  \end{enumerate}
\observatie Pentru acest criteriu, concluziile sunt inverse fata de criteriile precedente.
\end{tcolorbox}

\begin{tcolorbox}
\propozitie (Criteriul integral) Fie $f : [p, \infty) \to [0, \infty)$, $f$ descrescatoare. Atunci seria $\sum_{n \geq p} f(n)$ convergenta $\Leftrightarrow$ sirul $\left(\int_p^n f(x) dx \right)_{n \geq p}$ este convergent.
\end{tcolorbox}

\end{enumerate}
\section{Serii absolut convergente}
\begin{enumerate}

  \definitie Fie $(z_n)_{n \geq q}$ un sir de numere complexe. Seria $\sum_{n \geq q} z_n$ se numeste \term{absolut convergenta} daca seria $\sum_{n \geq q} \abs{z_n}$ este convergenta.
  \propozitie Orice serie absolut convergenta este si convergenta.

  \propozitie (Criteriul Dirichlet) Fie o serie $\sum_{n \geq q} a_n z_n$,  $a_n \in \mathbb{R}_+$,  $z_n \in \mathbb{C}, \forall n \geq q$. Daca $\exists N \geq q$ astfel incat: \\ $a_{n+1} \leq a_n, \forall n \geq N$ (sirul este descrescator dupa un punct) si $a_n \to 0$, apoi \\ $\exists M > 0$ a.i $\forall n \geq N$ avem $\abs{z_N + z_{N+1} + ... + z_{n}} \leq M$ \\
Atunci seria $\sum_{n \geq q} a_n z_n$ este convergenta.
  
\begin{tcolorbox}
\propozitie (Criteriul Leibniz) Fie seria $\sum_{n \geq q} (-1)^n a_n$. Daca $\exists N \in \mathbb{N} \geq q$ a.i $\forall n \geq N$ avem $a_{n+1} \leq a_n$ si $a_n \to 0$, atunci $\sum_{n \geq q} (-1)^n a_n$ convergenta.
\end{tcolorbox}

\begin{tcolorbox}
\propozitie (Criteriul Abel) Fie seria $\sum_{n \geq q} a_n z_n$, $a_n \in \mathbb{R}, z_n \in \mathbb{C}$. Daca $\exists N \geq q$ a.i $\forall n \geq N, a_{n+1} \leq a_n$ sau $a_{n+1} \geq a_n$ (sirul este monoton dupa un punct), $a_n$ convergent si $\sum_{n} z_n$ convergenta, atunci $\sum_{n \geq q} a_n z_n$ este convergenta.
\end{tcolorbox}


  

\end{enumerate}
\section{Siruri de functii cu valori complexe}
\begin{enumerate}

  \definitie Fie $(f_n)_{n \geq q}$, $f_n : A \to \mathbb{C}$, $A \neq \varnothing$ un sir de functii cu valori complexe si $f : A \to \mathbb{C}$. Spunem ca $f_n$ \term{converge simplu} (punctual) pe A la $f$ daca $\forall x \in A$, $\lim_{n \to \infty} f_n(x) = f(x)$.
\notatie Daca $f_n$ converge simplu pe A la $f$ se noteaza $f_n \xrightarrow[A]{s} f$.

\definitie Fie $(f_n)_{n \geq q}$, $f_n : A \to \mathbb{C}$, $A \neq \varnothing$ un sir de functii cu valori complexe si $f : A \to \mathbb{C}$. Spunem ca $f_n$ \term{converge uniform} pe A la $f$ daca $\forall \varepsilon > 0, \exists n_{\varepsilon} \geq q$ a.i $\forall n \geq n_{\varepsilon}, x \in A$ avem $\abs{f_n(x) - f(x)} < \varepsilon$.
\notatie Daca $f_n$ converge uniform pe A la $f$ se noteaza $f_n \xrightarrow[A]{u} f$.

\proprietate Daca $f_n \xrightarrow[A]{u} f$, atunci $f_n \xrightarrow[A]{s} f$.
\proprietate Daca $f_n \xrightarrow[A]{u} f$ si $B \subseteq A$ atunci $f_n \xrightarrow[B]{s} f$.
\proprietate Daca $A = A_1 \cup A_2$, $f_n \xrightarrow[A_1]{u} f$ si $f_n \xrightarrow[A_2]{u} f$, atunci $f_n \xrightarrow[A]{u} f$.
\proprietate Daca $f_n \xrightarrow[A]{u} f$, atunci $\abs{f_n - f} \xrightarrow[A]{u} 0$.
\proprietate $f_n \xrightarrow[A]{u} 0 \Leftrightarrow \abs{f_n} \xrightarrow[A]{u} 0$.
\proprietate Daca $f_n \xrightarrow[A]{u} f$, atunci $\abs{f_n} \xrightarrow[A]{u} \abs{f}$.
\proprietate Daca $f_n = \varphi, \forall n \geq q$, atunci $f_n \xrightarrow[A]{u} \varphi$.
\proprietate Daca $(z_n)_{n \geq q}$ este un sir de numere complexe, $z_n \to L$ si $f_n(x) = z_n, \forall n \geq q, x \in A$, atunci $f_n \xrightarrow[A]{u} L$.

\begin{tcolorbox}
\proprietate $f_n \xrightarrow[A]{u} f \Leftrightarrow \sup_{x \in A}  \abs{f_n(x) - f(x)} \to 0$.
\end{tcolorbox}

\proprietate Fie $f_n, f : I \to \mathbb{C}$, $I \subseteq \mathbb{R}$ interval, $n \geq q$ si $g : I \to \mathbb{C}$. Daca $f_n \xrightarrow[I]{s} f$, $f_n$ derivabile si $f_n' \xrightarrow[I]{u} g$, atunci $f$ e derivabila pe $I$ si $f' = g$. In plus, daca $I$ este marginit, atunci $f_n \xrightarrow[I]{u} f$.

\proprietate Daca $f_n \xrightarrow[A]{u} f$ si $g_n \xrightarrow[A]{u} g$, atunci $\alpha f_n + \beta g_n \xrightarrow[A]{u} \alpha f + \beta g$, $\forall \alpha, \beta \in \mathbb{C}$.

\proprietate Daca $f_n \xrightarrow[A]{u}f$, atunci $\alpha f_n \xrightarrow[A]{u} \alpha f$, $\forall \alpha \in \mathbb{C}$.

\proprietate Daca $f_{k_n}$ este un subsir al sirului $f_n \xrightarrow[A]{u} f$, atunci $f_{k_n} \xrightarrow[A]{u} f$.

  \definitie Fie $(f_n)_{n \geq q}, f_n : A \to \mathbb{C}$ un sir de functii cu $A \neq \varnothing$. Atunci $f_n$ se numeste \term{uniform Cauchy} pe $A$ daca $\forall \varepsilon > 0, \exists n_{\varepsilon} \geq q$ a.i $\forall n \geq n_{\varepsilon}, p \geq 1, x \in A$ avem $\abs{f_{n + p}(x) - f_n(x)} < \varepsilon$.


\end{enumerate}
\section{Serii de functii}
\begin{enumerate}

  \definitie Fie $f_n : A \to \mathbb{C}$, $A \neq \varnothing$, $n \geq q$. Seria $\sum_{n \geq q} f_n$ se numeste \term{convergenta simplu} pe $A$ daca $\forall x \in A$ seria $\sum_{n \geq q} f_n(x)$ este convergenta.

  \definitie Fie $f_n : A \to \mathbb{C}$, $A \neq \varnothing$, $n \geq q$. Seria $\sum_{n \geq q} f_n$ se numeste \term{convergenta uniform} pe $A$ daca sirul sumelor partiale asociat $S_n(x) = f_q(x) + f_{q+1}(x) + \dots + f_{n}(x)$ este uniform convergent pe $A$.
  \proprietate Daca $\sum_{n \geq q} f_n$ e convergenta uniform pe $A$, atunci $\sum_{n \geq q} \alpha f_n $ este uniform convergenta pe $A$, $\forall \alpha \in \mathbb{C}$.
  \proprietate Daca $\sum_{n \geq q} f_n$ si $\sum_{n \geq q} g_n$ sunt convergente uniform pe $A$, atunci $\sum_{n \geq q} (\alpha f_n + \beta g_n)$ este uniform convergenta pe $A$, $\forall \alpha, \beta \in \mathbb{C}$.
  \proprietate Fie $(z_n)_{n \geq q}$ un sir de numere complexe. Daca $\sum_{n \geq q} z_n$ este convergenta, atunci e si uniform convergenta pe orice multime nevida $A$.

  \propozitie (Criteriul necesar) Daca $\sum_{n \geq q} f_n$ e uniform convergenta pe $A$, atunci $f_n \xrightarrow[A]{u} 0$.
  
  \propozitie Daca $f_n$ este continua $\forall n \geq q$ si $\sum_{n \geq q} f_n$ e uniform convergenta pe intervalul $I$, atunci $S_n$ e continua pe $I$, $S_n \xrightarrow[I]{u} S$ si $\int_{a}^{b} S_n(x) \to \int_{a}^{b} S(x)$, $\forall [a, b] \subseteq I$.

\begin{tcolorbox}
  \teorema (Transfer de continuitate) Fie $f_n : I \to \mathbb{C}$, $I \subseteq \mathbb{R}$ interval, $f_n$ continua $\forall n \geq q$. Daca $\sum_{n \geq q} f_n$ e uniform convergenta pe $I$ atunci suma seriei: $\sum_{n = q}^{\infty} f_n$ este continua pe $I$.
\end{tcolorbox}

\begin{tcolorbox}
  \teorema (Integrare termen cu termen) Fie $f_n : [a, b] \to \mathbb{C}$, $f_n$ continua pe $[a, b]$, $\forall n \geq q$. Daca $\sum_{n \geq q} f_n$ este uniform convergenta pe $[a, b]$, atunci $\sum_{n \geq q} \left(\int_a^b f_n(x) dx \right)$ este convergenta si $\sum_{n = q}^{\infty} \left(\int_a^b f_n(x) dx \right) = \int_a^b \sum_{n = q}^{\infty} f_n(x)$.
\end{tcolorbox}

\begin{tcolorbox}
  \teorema (Derivare termen cu termen) Fie $f_n : I \to \mathbb{C}$, $I \subseteq \mathbb{R}$ interval, $f_n$ derivabila pe $I$, $\forall n \geq q$. Daca $\sum_{n \geq q}f_n$ e convergenta simplu pe $I$ si $\sum_{n \geq q} f'_n$ e convergenta uniform pe $I$, atunci $\sum_{n = q}^{\infty} f_n$ e derivabila si $\left(\sum_{n = q}^{\infty} f_n \right)' = \sum_{n = q}^{\infty} f_n'$.
\end{tcolorbox}



  \propozitie Fie $f_n : A \to \mathbb{C}$, $\forall n \geq q$ si $g_n : A \to \mathbb{R}$, $\forall n \geq N$, $N \geq q$. Daca $\abs{f_n(x)} \leq g_n(x)$, $\forall n \geq N, x \in A$, iar $\sum_{n \geq N} g_n$ e uniform convergenta pe $A$, atunci $\sum_{n \geq N} f_n$ e uniform convergenta pe $A$.

  \propozitie (Criteriul Cauchy pentru serii de functii). Fie $f_n : A \to \mathbb{C}$, $A \neq \varnothing$, $n \geq q$ un sir de functii. Atunci $\sum_{n \geq q} f_n$ uniform convergenta pe $A$ $\Leftrightarrow \forall \varepsilon > 0, \exists n_{\varepsilon} \geq q$ a.i $\forall n \geq n_{\varepsilon}, p \geq 1, x \in A$ avem $\abs{f_{n+1} + f_{n+2} + \dots + f_{n + p}} < \varepsilon$.

  \observatie Pentru siruri de functii cu valori complexe, definitiile pentru sir uniform Cauchy si sir uniform convergent sunt echivalente. (Atentie! Nu e valabil pentru orice sir de functii)

  
\begin{tcolorbox}
  \propozitie (Criteriul Weierstrass) Fie $f_n : A \to \mathbb{C}$, $A \neq \varnothing$, $n \geq q$. Daca $\exists (a_n)_{n \geq N}$ un sir de numere reale astfel incat $\abs{f_n(x)} \leq a_n, \forall n \geq N, x \in A$, iar $\sum_{n \geq N} a_n$ e convergenta, atunci $\sum_{n \geq q} f_n$ e uniform convergenta.
\end{tcolorbox}

  \propozitie (Criteriul Dirichlet pentru serii de functii). Fie $f_n : A \to \mathbb{R}$, $g_n : A \to \mathbb{C}$, $n \geq q$. Daca $f_{n+1}(x) \leq f_n(x)$, $\forall n \geq q, x \in A$, $f_n \xrightarrow[A]{u} 0$ si $\exists M > 0, r \geq q$ a.i $\forall n \geq r, x \in A$ avem $\abs{g_r(x) + g_{r+1}(x) + \dots + g_n(x)} \leq M$ (sirul sumelor partiale pentru seria $\sum_{n \geq q} g_n$ e egal marginit), atunci $\sum_{n \geq q} \left(f_n(x) \cdot g_n(x)\right)$ e uniform convergenta pe $A$.

\begin{tcolorbox}
  \propozitie (Criteriul Leibniz pentru serii de functii). Fie $f_n : A \to \mathbb{R}$, $n \geq q$ un sir de functii, $A \neq \varnothing$. Daca $\exists N \geq q$ a.i $\forall n \geq N, x \in A$ avem $f_{n+1}(x) \leq f_n(x)$ si $f_n \xrightarrow[A]{u} 0$, atunci $\sum_{n \geq q} (-1)^n f_n(x)$ e uniform convergenta pe $A$.
\end{tcolorbox}

\begin{tcolorbox}
  \propozitie (Criteriul Abel pentru serii de functii). Fie $f_n : A \to \mathbb{R}$, $g_n : A \to \mathbb{C}$, $n \geq q$.  Daca $f_{n+1}(x) \leq f_n(x)$ (sau $f_{n+1}(x) \geq f_n(x)$), $\forall n \geq q$, $\exists M > 0, N \geq q$ a.i $\forall n \geq q, x \in A$ avem $\abs{f_n(x)} \leq M$ si $\sum_{n \geq q} g_n$ e uniform convergenta pe $A$, atunci $\sum_{n \geq q} (f_n \cdot g_n)$ e uniform convergenta pe $A$.
\end{tcolorbox}
  



\end{enumerate}
\section{Serii de puteri}
\begin{enumerate}
  \definitie Se numeste o \term{serie de puteri} o serie de forma $\sum_{n \geq 0} a_n z^n$, unde $(a_n)_{n \geq 0}$ e un sir de numere complexe si $z \in \mathbb{C}$.
  \observatie Pentru $z = 0$ seria de puteri este convergenta si suma ei este egala cu $a_0$.
  
  \definitie Fie $\sum_{n \geq 0} a_n z^n$ o serie de puteri si $R \xlongequal{\text{not}} \sup \left\{r \in \mathbb{R}_+ \middle\vert \sum_{n \geq 0} \abs{a_n} r^n \text{ convergenta } \right\}$.

\begin{tcolorbox}
  \propozitie Fie $R$ raza de convergenta a seriei $\sum_{n \geq 0} a_n z^n$. Atunci: 
  \begin{enumerate}
    \item Daca $\abs{z} < R$, atunci $\sum_{n \geq 0} a_n z^n$ e absolut convergenta.
    \item Daca $\abs{z} > R$, atunci $\sum_{n \geq 0} a_n z^n$ e divergenta.
  \end{enumerate}
  \observatie Daca $R = \infty$, seria e convergenta pentru orice $z$.
  \observatie Daca $R = 0$, seria e divergenta pentru orice $z$ nenul.
\end{tcolorbox}

\proprietate Fie $R$ raza de convergenta a seriei $\sum_{n \geq 0} a_n x^n$, $x \in \mathbb{R}, R > 0$. Atunci, $\forall r \in (0, R)$, $\sum_{n \geq 0} a_n x^n$ este uniform convergenta pe $[-r , r]$.

  \propozitie Fie $R$ raza de convergenta a seriei $\sum_{n \geq 0} a_n z^n$. Atunci: 
  \begin{enumerate}
    \item Daca $\sum_{n \geq 0} a_n z^n$ e convergenta, atunci $\abs{z} \leq R$.
    \item Daca $\sum_{n \geq 0} a_n z^n$ e divergenta, atunci $\abs{z} \geq R$.
  \end{enumerate}

\begin{tcolorbox}
  \propozitie Fie $R$ raza de convergenta a seriei $\sum_{n \geq 0} a_n z^n$. Daca $\exists \lim_{n \to \infty} \sqrt[n]{\abs{a_n}} \xlongequal{\text{not}} L \in [0, \infty]$, atunci $R = \frac{1}{L}$.

\corolar Fie $R$ raza de convergenta a seriei $\sum_{n \geq 0} a_n z^n$. Daca $\exists N \in \mathbb{N}$ astfel incat \\ $a_n \neq 0, \forall n \geq N$ si $\exists \lim_{n \to \infty} \frac{\abs{a_{n+1}}}{\abs{a_n}} \xlongequal{\text{not}} L \in [0, \infty]$, atunci $R = \frac{1}{L}$.
\end{tcolorbox}


\propozitie 
  \begin{enumerate}
    \item Seriile $\sum_{n \geq 0} a_n x^n$ si $\sum_{n \geq 1} n a_n x^{n-1}$ au aceeasi raza de convergenta.
    \item Seriile $\sum_{n \geq 0} a_n x^n$ si $\sum_{n \geq 0} \frac{a_n}{n+1} x^{n+1}$ au aceeasi raza de convergenta. 
  \end{enumerate}

\begin{tcolorbox}
  \propozitie Fie $R$ raza de convergenta a seriei $\sum_{n \geq 0} a_n x^n$, $R > 0$. Fie $S : (-R, R) \to \mathbb{C}$, $S(x) = \sum_{n = 0}^{\infty} a_n x^n$ suma seriei. Atunci $S$ este infinit derivabila pe $(-R, R)$ si se poate deriva termen cu termen de oricate ori pe $(-R, R)$.
  \observatie $S'(x) = \sum_{n = 1}^{\infty} n a_n x^{n-1}$, $\forall x \in (-R, R)$.
\end{tcolorbox}

  \propozitie Fie $R$ raza de convergenta a seriei $\sum_{n \geq 0} a_n x^n$, $R > 0$. Fie $S : (-R, R) \to \mathbb{C}$, $S(x) = \sum_{n = 0}^{\infty} a_n x^n$ suma seriei si $F : (-R, R) \to \mathbb{C}$ o primitiva a lui $S$. Atunci $\exists c \in \mathbb{C}$ astfel incat $\sum_{n = 0}^{\infty} \frac{a_n}{n + 1} x^{n+1} = F(x) + c$, $\forall x \in (-R, R)$. 

\begin{tcolorbox}
  \propozitie Fie $R$ raza de convergenta a seriei $\sum_{n \geq 0} a_n x^n$, $R > 0$. \\ Fie $S : (-R, R) \to \mathbb{C}$ suma seriei. Atunci $\forall [a, b] \subseteq (-R, R)$, $S(x)$ se poate integra pe $[a, b]$ si $\int_a^b S(x) dx = \int_a^b \left(\sum_{n = 0}^{\infty} a_n x^n \right)dx = \sum_{n=0}^{\infty}\int_a^b a_n x^n dx$.
\end{tcolorbox}

\newpage
  \propozitie Fie $R$ raza de convergenta a seriei $\sum_{n \geq 0} a_n x^n$, $R \in (0, \infty)$ si $S$ suma seriei. Daca $\sum_{n \geq 0} a_n R^n$ e convergenta, atunci:
  \begin{enumerate}[nosep]
    \item $\sum_{n \geq 0} a_n x^n$ uniform convergenta pe $[0, R]$.
    \item $S$ continua pe $[0, R]$.
    \item $\sum_{n = 0}^{\infty} a_n R^n = \lim_{\substack{x \to R \\ x < R}} S(x)$.
    \item $\int_0^R S(x) dx = \sum_{n = 0}^{\infty} \frac{a_n}{n+1} R^{n+1}$.
  \end{enumerate}

  \propozitie Fie $\sum_{n \geq 0} a_n$ convergenta, $a_n \in \mathbb{C}, \forall n \geq 0$. Atunci raza de convergenta a seriei este $R \geq 1$ si $\sum_{n = 0}^{\infty} a_n = \lim_{\substack{ x \to 1 \\ x < 1}} \sum_{n = 0}^{\infty} a_n x^n$.

\begin{tcolorbox}
  \propozitie Fie $(a_n)_{n \geq 0}$ sir de numere reale. Daca $\exists p \in \mathbb{N}$ astfel incat $a_{n+1} \leq a_n$, $\forall n \geq p$ si $a_n \to 0$, atunci $\sum_{n \geq p} (-1)^n a_n$ e convergenta. \\ Fie $S$ suma acestei serii si $(S_n)_{n \geq p}$ sirul de sume partiale asociat. Atunci $\abs{S_n - S} \leq a_{n+1}$ si:
  \begin{enumerate}
  \item Daca $S_n = \dots + a_n$, atunci $S_n \geq S$.
  \item Daca $S_n = \dots - a_n$, atunci $S_n \leq S$.
  \end{enumerate}

  \observatie Aceasta propozitie poate fi folosita pentru aproximarea numerica a unor expresii. Mai intai se exprima ca suma unei astfel de serii, apoi se poate incadra intre doua sume partiale consecutive.

\end{tcolorbox}

\propozitie Raza de convergenta a seriei $\sum_{n \geq 0} \frac{z^n}{n!}$ este $\infty$.

\begin{tcolorbox}
\propozitie Urmatoarele functii pot fi exprimate ca sumele unor serii de puteri:
  \begin{enumerate}[nosep]
    \item $\ln(1+x) = \sum_{n=1}^{\infty} (-1)^{n-1} \cdot \frac{x^n}{n}$, $\forall x \in (-1, 1]$.
    \item $\arctan x = \sum_{n = 0}^{\infty} (-1)^n \cdot \frac{x^{2n + 1}}{2n + 1}$, $\forall x \in [-1, 1]$.
    \item $e^z = \sum_{n = 0}^{\infty} \frac{z^n}{n!}$, $\forall z \in \mathbb{C}$.
    \item $\cos(z) = \frac{e^{iz} + e^{-iz}}{2} = \sum_{n=0}^{\infty} (-1)^n \cdot \frac{z^{2n}}{(2n)!}$
    \item $\sin(z) = \frac{e^{iz} - e^{-iz}}{2i} = \sum_{n=0}^{\infty} (-1)^n \cdot \frac{z^{2n+1}}{(2n+1)!}$
    \item $\cosh(z) \xlongequal{\text{def}} \frac{e^z + e^{-z}}{2}$, $\forall z \in \mathbb{C}$.
    \item $\sinh(z) \xlongequal{\text{def}} \frac{e^z - e^{-z}}{2}$, $\forall z \in \mathbb{C}$.

  \end{enumerate}
  \observatie Extinderea pe $\mathbb{C}$ a functiilor $e^x$, $\cos(x)$, $\sin(x)$, etc. este definita prin aceste serii.
  \observatie $e^{x + iy} = e^x(\cos(y) + i\sin(y))$.
\end{tcolorbox}


\end{enumerate}
\section{Formula lui Taylor}
\begin{enumerate}
\begin{tcolorbox}
  \teorema (Formula lui Taylor) Fie $f : I \to \mathbb{R}$, $I \subseteq \mathbb{R}$ interval, $f$ de $n+1$ ori derivabila, $n \in \mathbb{N}$ si $a \in I$. Atunci $\forall x \in I$, $x \neq a$, $\exists c \in [a, x]$ (sau $[x, a]$) astfel incat: 
  $$f(x) = \sum_{k = 0}^{n} \frac{f^{(k)}(a)}{k!} (x-a)^k + \frac{f^{(n+1)}(c)}{(n+1)!} (x-a)^{n+1}$$

  \notatie $T_{n, a} \xlongequal{\text{not}} \sum_{k = 0}^{n} \frac{f^{(k)}(a)}{k!} (x-a)^k$ se numeste \term{polinom Taylor de ordin $n$}.
  \observatie Pentru $a=0$ se obtine formula lui Maclaurin.
  \observatie $T_{n, a} = f(a) + \frac{f'(a)}{1!} (x-a) + \frac{f''(a)}{2!} (x-a)^2 + \dots + \frac{f^{(n)}(a)}{n!} (x-a)^n$.
\end{tcolorbox}
\corolar Fie $P \in \mathbb{R}[X]$ cu $\deg P \leq n$. Atunci $P(X) = \sum_{k = 0}^{n} \frac{P^{(k)}(a)}{k!} (X-a)^k$.

\teorema (Mertens) Fie $\sum_{n \geq 0} A_n$ si $\sum_{n \geq 0} B_n$ doua serii de numere complexe cu $\sum_{n \geq 0} A_n$ absolut convergenta si $\sum_{n \geq 0} B_n$ convergenta. Fie $C_n = A_0 B_n + A_1 B_{n-1} + \dots + A_n B_0 = \sum_{k=0}^{n} A_k B_{n-k}$. Atunci $\sum_{n \geq 0} C_n$ convergenta si $\sum_{n=0}^{\infty} C_n = \left( \sum_{n = 0}^{\infty} A_n \right) \cdot \left( \sum_{n = 0}^{\infty} B_n \right)$


\end{enumerate}
\end{document}
